% GENERATED FILE. Edit canonical lessons, then run npm run generate:books.
% canonical-source-hash: fnv1a64:3ee67aba236a52b2
% canonical-lessons: GE-C437-bummeln, GE-C437-schlendern, GE-C437-erkundigen, GE-C437-auskunft, GE-C437-auskennen, GE-R437-review-430-437

\chapter{Durch die Stadt}
\label{ch:ge-durch-die-stadt}
\hlchaptermodality{437}

\begin{chapteropening}
\textbf{By the end of this chapter:} \emph{I can wander round a town and ask the way: strolling, sauntering, asking about something, information and knowing my way round.}

Complete the last lesson of chapter 437: I can wander round a town and ask the way: strolling, sauntering, asking about something, information and knowing my way round.
\end{chapteropening}

\section[bummeln]{bummeln — to stroll, to wander around; to dawdle}

\label{lesson:GE-C437-bummeln}



\begin{quote}
Before the new one: say the German for the one-way street, then the German for the passer-by.
\end{quote}

\subsection*{You'll want to know: bummeln}
\textbf{bummeln} — \textquotedblleft{}to stroll, to wander around; to dawdle\textquotedblright{}.

\begin{quote}
Am Samstag bummeln wir durch die Fußgängerzone.
\end{quote}

\begin{quote}
\emph{On Saturday we stroll through the pedestrian zone.}
\end{quote}

Perfekt with \emph{sein} for moving: \emph{wir sind durch die Stadt gebummelt}. \emph{Der Stadtbummel} — a stroll round town.

\begin{cousinweb}
The origin of \emph{bummeln} is not settled. It is often explained as a sound word for something swinging slowly to and fro, like a bell; EtymDB lists the sense 'to loaf'.
\end{cousinweb}

\begin{grammarlens}[title={What you've built}]
One more word for telling, explaining and describing at B1.
\end{grammarlens}

\subsection*{Your turn}
\begin{itemize}
\raggedright
  \item \emph{Say it:} \emph{bummeln}
  \item \emph{Say it:} \emph{bummeln}, once more
  \item \emph{Say it:} a sentence of your own with \emph{bummeln}
  \item \emph{Recall:} say the German for the one-way street, then the German for the passer-by, then say \emph{bummeln} again
\end{itemize}

\begin{tcolorbox}[breakable,colback=teal!4,colframe=teal!35!black,title={Before you move on}]
What does bummeln mean? (\textquotedblleft{}To stroll, to wander around; to dawdle\textquotedblright{}.) Say it once more.
\end{tcolorbox}
\section[schlendern]{schlendern — to saunter, to amble}

\label{lesson:GE-C437-schlendern}



\begin{quote}
Before the new one: say the German for the passer-by, then the German for to stroll.
\end{quote}

\subsection*{You'll want to know: schlendern}
\textbf{schlendern} — \textquotedblleft{}to saunter, to amble\textquotedblright{}.

\begin{quote}
Nach dem Essen schlendert er langsam zum Hotel.
\end{quote}

\begin{quote}
\emph{After the meal he saunters slowly back to the hotel.}
\end{quote}

Perfekt with \emph{sein}: \emph{wir sind am Ufer entlanggeschlendert}.

\begin{cousinweb}
\emph{schlendern} came into the standard language from Low German. EtymDB links it with Middle High German \emph{lenderen} 'to loiter'; the older history is not settled.
\end{cousinweb}

\begin{grammarlens}[title={What you've built}]
One more word for telling, explaining and describing at B1.
\end{grammarlens}

\subsection*{Your turn}
\begin{itemize}
\raggedright
  \item \emph{Say it:} \emph{schlendern}
  \item \emph{Say it:} \emph{schlendern}, once more
  \item \emph{Say it:} a sentence of your own with \emph{schlendern}
  \item \emph{Recall:} say the German for the passer-by, then the German for to stroll, then say \emph{schlendern} again
\end{itemize}

\begin{tcolorbox}[breakable,colback=teal!4,colframe=teal!35!black,title={Before you move on}]
What does schlendern mean? (\textquotedblleft{}To saunter, to amble\textquotedblright{}.) Say it once more.
\end{tcolorbox}
\section[erkundigen]{erkundigen — to ask about, to enquire (mostly reflexive: sich erkundigen)}

\label{lesson:GE-C437-erkundigen}



\begin{quote}
Before the new one: say the German for to stroll, then the German for to saunter.
\end{quote}

\subsection*{You'll want to know: erkundigen}
\textbf{erkundigen} — \textquotedblleft{}to ask about, to enquire (mostly reflexive: sich erkundigen)\textquotedblright{}.

\begin{quote}
Sie erkundigt sich nach dem Weg zum Bahnhof.
\end{quote}

\begin{quote}
\emph{She asks the way to the station.}
\end{quote}

Reflexive, with \emph{nach}: \emph{ich erkundige mich nach den Öffnungszeiten}.

\begin{cousinweb}
\emph{erkundigen} is built from \emph{er-} and \emph{kundig} 'well informed', from the old adjective \emph{kund} 'known', an inherited Germanic word related to \emph{kennen} 'to know'.
\end{cousinweb}

\begin{grammarlens}[title={What you've built}]
One more word for telling, explaining and describing at B1.
\end{grammarlens}

\subsection*{Your turn}
\begin{itemize}
\raggedright
  \item \emph{Say it:} \emph{erkundigen}
  \item \emph{Say it:} \emph{erkundigen}, once more
  \item \emph{Say it:} a sentence of your own with \emph{erkundigen}
  \item \emph{Recall:} say the German for to stroll, then the German for to saunter, then say \emph{erkundigen} again
\end{itemize}

\begin{tcolorbox}[breakable,colback=teal!4,colframe=teal!35!black,title={Before you move on}]
What does erkundigen mean? (\textquotedblleft{}To ask about, to enquire (mostly reflexive: sich erkundigen)\textquotedblright{}.) Say it once more.
\end{tcolorbox}
\section[die Auskunft]{die Auskunft — the information; the information desk}

\label{lesson:GE-C437-auskunft}



\begin{quote}
Before the new one: say the German for to saunter, then the German for to ask about.
\end{quote}

\subsection*{You'll want to know: die Auskunft}
\textbf{die Auskunft} — \textquotedblleft{}the information; the information desk\textquotedblright{}.

\begin{quote}
Am Bahnhof gibt uns eine Mitarbeiterin Auskunft.
\end{quote}

\begin{quote}
\emph{At the station a member of staff gives us information.}
\end{quote}

Plural: \emph{die Auskünfte}. \emph{Die Auskunft} is also the information desk itself.

\begin{cousinweb}
\emph{Auskunft} is usually traced to Middle High German \emph{ūzkunft} 'a way out', built from \emph{aus} 'out' and \emph{Kunft}, from \emph{kommen} 'to come'; the sense 'information' came later. EtymDB does not list the word, so the line is hedged.
\end{cousinweb}

\begin{grammarlens}[title={What you've built}]
One more word for telling, explaining and describing at B1.
\end{grammarlens}

\subsection*{Your turn}
\begin{itemize}
\raggedright
  \item \emph{Say it:} \emph{die Auskunft}
  \item \emph{Say it:} \emph{die Auskunft}, once more
  \item \emph{Say it:} a sentence of your own with \emph{die Auskunft}
  \item \emph{Recall:} say the German for to saunter, then the German for to ask about, then say \emph{die Auskunft} again
\end{itemize}

\begin{tcolorbox}[breakable,colback=teal!4,colframe=teal!35!black,title={Before you move on}]
What does die Auskunft mean? (\textquotedblleft{}The information; the information desk\textquotedblright{}.) Say it once more.
\end{tcolorbox}
\section[auskennen]{auskennen — to know one's way around (mostly reflexive: sich auskennen)}

\label{lesson:GE-C437-auskennen}



\begin{quote}
Before the new one: say the German for to ask about, then the German for the information.
\end{quote}

\subsection*{You'll want to know: auskennen}
\textbf{auskennen} — \textquotedblleft{}to know one's way around (mostly reflexive: sich auskennen)\textquotedblright{}.

\begin{quote}
In dieser Stadt kenne ich mich gut aus.
\end{quote}

\begin{quote}
\emph{I know my way around this town well.}
\end{quote}

Reflexive and separable: \emph{ich kenne mich aus}, \emph{ich habe mich ausgekannt}.

\begin{cousinweb}
\emph{auskennen} is built from \emph{aus} 'out, through' and \emph{kennen} 'to know', inherited from Old High German \emph{kennan} and Proto-Germanic \emph{kannijaną} (EtymDB).
\end{cousinweb}

\begin{grammarlens}[title={What you've built}]
One more word for telling, explaining and describing at B1.
\end{grammarlens}

\subsection*{Your turn}
\begin{itemize}
\raggedright
  \item \emph{Say it:} \emph{auskennen}
  \item \emph{Say it:} \emph{auskennen}, once more
  \item \emph{Say it:} a sentence of your own with \emph{auskennen}
  \item \emph{Recall:} say the German for to ask about, then the German for the information, then say \emph{auskennen} again
\end{itemize}

\begin{tcolorbox}[breakable,colback=teal!4,colframe=teal!35!black,title={Before you move on}]
What does auskennen mean? (\textquotedblleft{}To know one's way around (mostly reflexive: sich auskennen)\textquotedblright{}.) Say it once more.
\end{tcolorbox}
\section[(dialogue)]{Words from chapters 430–437 — a review}

\label{lesson:GE-R437-review-430-437}



\begin{quote}
Say the German for the bedbug and the cockroach.
\end{quote}

\subsection*{Your turn}
Say these aloud:

\begin{itemize}
\raggedright
  \item the bedbug, the cockroach, the mosquito, to bite, to crawl
  \item the complaint, to complain, to take back, to grumble, dissatisfied
  \item the manager, the inconvenience, the deposit, to compensate, the compensation
  \item to get one's bearings, the compass, the point of the compass, to find one's way around, to get lost
  \item the signpost, the sign, the arrow, the turning, to branch off
  \item the shortcut, to take a shortcut, the detour, to turn back, to get lost driving
  \item the lane, the dead end, the avenue, the one-way street, the passer-by
  \item to stroll, to saunter, to ask about, the information, to know one's way around
\end{itemize}

\begin{tcolorbox}[breakable,colback=teal!4,colframe=teal!35!black,title={Before you move on}]
The bedbug? (\textbf{die Wanze}.) The cockroach? (\textbf{die Kakerlake}.) The mosquito? (\textbf{die Mücke}.) To bite? (\textbf{stechen}.) To crawl? (\textbf{krabbeln}.) The complaint? (\textbf{die Beschwerde}.) To complain? (\textbf{beschweren}.) To take back? (\textbf{reklamieren}.) To grumble? (\textbf{meckern}.) Dissatisfied? (\textbf{unzufrieden}.) The manager? (\textbf{der Geschäftsführer}.) The inconvenience? (\textbf{die Unannehmlichkeit}.) The deposit? (\textbf{die Anzahlung}.) To compensate? (\textbf{entschädigen}.) The compensation? (\textbf{die Entschädigung}.) To get one's bearings? (\textbf{orientieren}.) The compass? (\textbf{der Kompass}.) The point of the compass? (\textbf{die Himmelsrichtung}.) To find one's way around? (\textbf{zurechtfinden}.) To get lost? (\textbf{verirren}.) The signpost? (\textbf{der Wegweiser}.) The sign? (\textbf{das Schild}.) The arrow? (\textbf{der Pfeil}.) The turning? (\textbf{die Abzweigung}.) To branch off? (\textbf{abzweigen}.) The shortcut? (\textbf{die Abkürzung}.) To take a shortcut? (\textbf{abkürzen}.) The detour? (\textbf{der Umweg}.) To turn back? (\textbf{umkehren}.) To get lost driving? (\textbf{verfahren}.) The lane? (\textbf{die Gasse}.) The dead end? (\textbf{die Sackgasse}.) The avenue? (\textbf{die Allee}.) The one-way street? (\textbf{die Einbahnstraße}.) The passer-by? (\textbf{der Passant}.) To stroll? (\textbf{bummeln}.) To saunter? (\textbf{schlendern}.) To ask about? (\textbf{erkundigen}.) The information? (\textbf{die Auskunft}.) To know one's way around? (\textbf{auskennen}.)
\end{tcolorbox}
